% s-369-p369-m002; original PDF 369, 370
\Aphorism{s-369-p369-m002}{411}
Um den Durchseihungsprocess, vermöge dessen aus einem transcendenten Zufallsspiel eine empirische Causalwelt herausfliessen kann, mit einem wissenschaftlichen Namen zu taufen, nennen wir ihn den \Sperr{Process der indirecten Auslese} — in dem Sinne, wie die Biologie von indirecter Auslese spricht. Das will sagen, dass empirische Gesetzlichkeit, ebenso wie biologische Zweckmässigkeit, nicht von vornherein als einziger Fall verwirklicht wurde, sondern nur schliesslich als einziger übrig bleibt, vermöge einer selectiven Vorrichtung, die alle davon abweichenden Fälle beseitigt. Diese selective Vorrichtung, diese ins Erkenntnisstheoretische übersetzte Zuchtwahl heisst eben Bewusstsein — aber ihre Function ist nicht zeitlich, sondern ontologisch zu denken: das Herauskrystallisiren des Kosmos aus dem Chaos vollzieht sich nicht im Sinne eines Verlaufs, sondern etwa einer rein mathematischen Elimination. Sehen wir beispielsweise ein Lebewesen als Atomschwarm an, dessen Theilchen nur bei einer ganz bestimmten Anordnung das bewusste Lebewesen constituiren, sonst aber eine \Foreign{rudis indigestaque moles} vorstellen; denken wir uns dann die Theilchen beliebig den Weltraum durchfliegen und sich nur gelegentlich und zufällig zu jener Anordnung zusammenfinden — so werden die Zeiten, da die Theilchen ungeordnet herumirren, in dem Bewusstsein jenes Lebewesens nicht vertreten sein; es wird, sowie es sich stets als bewusst vorfindet, auch stets die sein Bewusstsein bedingende Anordnung der Atome realisirt vorfinden. Während für einen objectiven Weltbetrachter die Wahrscheinlichkeit jener Anordnung eine verschwindend geringe sein kann, ist sie für das Lebewesen selbst subjective Gewissheit; das Bewusstsein schneidet von selbst aus der Fläche aller Möglichkeiten die Linie heraus, längs deren allein die Bedingungen dieses Bewusstseins erfüllt sind. So ist jede Art Bewusstsein, welche an Bedingungen geknüpft ist — und dies scheint die uns allein vorstellbare Art zu sein — ein Sieb, ein Eliminationsapparat, eine Zwangsgleichung für sonst unabhängige Variable, eine Gesetzlichkeit innerhalb eines Zufallsspiels, ein Kosmos mitten im Chaos — wohlgemerkt, immer nur „subjectiv“, nur für den Träger des Bewusstseins, für Niemanden sonst! Danach wäre die Gesetzlichkeit der Welt einfach darauf zurückgeführt, dass die zahllosen \Sperr{abweichenden Fälle nicht registrirt} werden, eine Quelle, aus der ja jede Art Aberglauben Zufluss findet, der populäre, der wissenschaftliche, der philosophische! Und die allerkürzeste Formel der Erkenntnisstheorie hätte jener Knabe in einer Björnson’schen Novelle ausgesprochen, der seinem Spielgefährten die Vexirfrage stellt: „Hast du den Hahn in deinem Bilderbuch gesehen? Er hat eine Menge Hühner um sich, wenn das Buch zu ist: hast du das auch gesehn?“
\par

